% Included by sections/03-dataset.tex as the opening section of the dataset chapter
% (the acquisition is part of describing the dataset, not a chapter of its own).
\section{Data acquisition}
\label{sec:acquisition}

\subsection{Source: the FFVL Coupe F\'ed\'erale de Distance}

The data come from the \emph{Coupe F\'ed\'erale de Distance} (CFD), the cross-country
competition run by the French free-flight federation (FFVL). The federation operates two
parallel competitions on separate sites: one for paragliders
(\texttt{parapente.ffvl.fr}) and one for hang gliders (\texttt{delta.ffvl.fr}). Pilots
upload their flights; each entry carries metadata---\rev{among them} pilot, date, take-off
and landing sites, scored distance, duration, and glider\rev{, with the full set of fields
listed in Sec.~\ref{sec:catalog}}---and, in most cases, the raw GPS tracklog in IGC format
(the standard flight-recorder format, Sec.~\ref{sec:igcformat}).
A ``season'' is identified by its starting year,
e.g.\ the year~1999 denotes the 1999--2000 season; the paraglider archive starts in
1999--2000 and the hang-glider archive in 2001--2002.

\subsection{Collected data}

For every season the pipeline retrieves two things:
\begin{itemize}
  \item the \textbf{season XML export}, archived verbatim as provenance
        (\texttt{raw\_xml/<year>.xml}); it is the authoritative list of flights and links;
  \item the \textbf{IGC tracklog files} referenced by the XML, one file per flight that
        has a track.
\end{itemize}
From these, a derived \textbf{catalog} is built (\texttt{catalog.csv}, one row per flight,
\rev{its columns detailed in Sec.~\ref{sec:catalog}}) together with a per-season index
(\texttt{seasons\_index.csv}). The raw data (tens of gigabytes) live on an external disk
and are never committed to the repository. \flow{The catalog is itself rebuildable from
the archived XMLs and likewise uncommitted (Sec.~\ref{sec:catalog}).}

\subsection{Pipeline and tooling}

Acquisition is performed by a dedicated, installable Python package, exposed through two
command-line tools, one per discipline, sharing the same subcommands for fetching,
downloading, cataloguing, and verifying. The design favours robustness for a
long-running, large job: downloads are resumable and safe to interrupt, a response is
accepted only once validated as real IGC content, and download requests are rate-limited
so as not to load the server \impldetails{impl:acqpipeline}.

\subsection{Reproducibility}

The destination directory is not hard-coded: it is set per machine through an environment
variable or a configuration file, of which the repository ships only a placeholder, so no
machine-specific path is committed. Archiving the season
XMLs verbatim keeps the full provenance \impldetails{impl:acqrepro}.

\subsection{Integrity}

Because the data sit on an external drive, integrity is re-checked independently after
download, rather than trusted at download time \impldetails{impl:acqintegrity}. Both
FFVL collections have been verified in full: all \num{\StatParaDownloaded} paraglider and
\num{\StatHangDownloaded} hang-glider downloaded tracklogs pass validation against the
IGC standard (Sec.~\ref{sec:igcformat}). \flow{These are the tracklogs actually fetched; they cover a subset of the flights that
carry a track at all (Sec.~\ref{sec:coverage}).}

\subsection{Further data sources}
The FFVL CFD is the primary source of the present study, but the dataset is not assumed to
be final. The paraglider data (\num{\StatParaTotalFlights} flights,
\num{\StatParaWithIgc} tracklogs) are the main analysis target; the hang-glider data
(\num{\StatHangTotalFlights} flights, \num{\StatHangWithIgc} tracklogs) are collected in
parallel and enable the cross-type comparison of Sec.~\ref{sec:glider}. Additional
public soaring repositories are
being pursued to enlarge and diversify the sample and, above all, to add the third glider
category, \emph{sailplanes} (Sec.~\ref{sec:glider}). The most complete sailplane source
is \emph{WeGlide}, which exposes flight metadata and IGC tracks through an API. The
default access is rate-limited (of order sixty requests per day), which rules out a bulk
export of the archive; a request for research access with a higher quota has been sent
and is pending. Paraglider and hang-glider coverage may likewise be broadened through
further portals such as XContest. \rev{A comparable research-access request has been sent
to XContest; as of 20 July 2026 neither WeGlide nor XContest has replied.} All of these can
be acquired with the same approach\rev{, each producing a catalog in the same format}.
